% !PRINTED 535
\weberpage{535}

We have already seen above how, from one basis of $\mu$, all the
others can be derived.

We now understand by $\alpha_1, \alpha_2, \ldots, \alpha_n$ an
arbitrary basis of $\mu$, and represent the functional $\mu$,
according to \S.\,142, 4., as a quotient of two integral functions
whose denominator is a rational unit. The coefficients of the
numerator are then integers divisible by $\mu$ (\S.\,143, 2.), and can
therefore be represented in the form (2).

If we gather this representation together suitably, there thus results
for $\mu$ an expression
\weq{8}{
\mu = \alpha_1 u_1 + \alpha_2 u_2 + \cdots + \alpha_n u_n,
}
in which the $u_1, u_2, \ldots, u_n$ are \gesperrt{integral rational
functionals}.

With this we wish to compare the linear form
\weq{9}{
\lambda = \alpha_1 t_1 + \alpha_2 t_2 + \cdots + \alpha_n t_n
}
in which $t_1, t_2, \ldots, t_n$ are variables. This linear form is
(by \S.\,139) the greatest common divisor of the numbers
$\alpha_1, \alpha_2, \ldots, \alpha_n$, and is divisible by $\mu$,
because the numbers $\alpha_1, \alpha_2, \ldots, \alpha_n$ are
divisible by $\mu$. On the other hand, however, as the representation
(8) shows, $\mu$ is also divisible by $\lambda$, and consequently the
two functionals $\mu$ and $\lambda$ are associated with one another.

\gesperrt{The functional $\lambda$ we shall call a basis form of the
functional $\mu$; $\lambda$ is then at the same time a basis form of
all functionals associated with $\mu$.}

A basis form $\lambda$ of $\mu$ has the property that from
it one obtains all the integers in $\Omega$ divisible by $\mu$ (or by
$\lambda$) when one sets integral rational numbers for the
variables.

The basis $\alpha_1, \alpha_2, \ldots, \alpha_n$ stands to the
functional $\mu$ in a relation similar to that in which the basis
$\omega_1, \omega_2, \ldots, \omega_n$ of the system $\mathfrak{o}$ of all
integers in $\Omega$ stands to the units. In fact the linear form
\weq{10}{
\tau = \omega_1 t_1 + \omega_2 t_2 + \cdots + \omega_n t_n
}
is a unit; for it is the greatest common divisor of
$\omega_1, \omega_2, \ldots, \omega_n$, and must therefore, as the
formula \S.\,145, (8) shows, be a divisor of 1, and hence a unit.

\gesperrt{Accordingly we shall call the functional $\tau$ a basis form of
$\mathfrak{o}$.}
